\documentclass[11pt,a4paper]{article}

% --- Encoding ---
\usepackage[utf8]{inputenc}
\usepackage[T1]{fontenc}
\usepackage[spanish,es-noshorthands]{babel}
\usepackage{lmodern}

% --- Layout ---
\usepackage[margin=2.5cm]{geometry}
\usepackage{parskip}
\setlength{\parindent}{0pt}
\setlength{\parskip}{0.5em plus 0.2em minus 0.1em}

% --- Tables and graphics ---
\usepackage{booktabs}
\usepackage{array}
\usepackage{enumitem}
\usepackage{xcolor}
\usepackage{graphicx}
\usepackage{tikz}
\usetikzlibrary{positioning,shapes.geometric,arrows.meta,fit,backgrounds,calc,decorations.pathreplacing}
\usepackage{caption}
\usepackage{amsmath,amssymb}

% --- Code listings ---
\usepackage{listings}
\definecolor{codegray}{rgb}{0.96,0.96,0.96}
\definecolor{codeborder}{rgb}{0.85,0.85,0.85}
\definecolor{codecomment}{rgb}{0.30,0.55,0.30}
\definecolor{codekeyword}{rgb}{0.10,0.30,0.70}
\definecolor{codestring}{rgb}{0.65,0.15,0.10}

\lstdefinestyle{cstyle}{
    language=C,
    backgroundcolor=\color{codegray},
    basicstyle=\ttfamily\footnotesize,
    breaklines=true,
    frame=single,
    rulecolor=\color{codeborder},
    showstringspaces=false,
    keywordstyle=\color{codekeyword}\bfseries,
    commentstyle=\color{codecomment}\itshape,
    stringstyle=\color{codestring},
    aboveskip=1em,
    belowskip=1em,
}

\lstdefinestyle{shellstyle}{
    language=bash,
    backgroundcolor=\color{codegray},
    basicstyle=\ttfamily\footnotesize,
    breaklines=true,
    frame=single,
    rulecolor=\color{codeborder},
    showstringspaces=false,
    keywordstyle=\color{codekeyword}\bfseries,
    commentstyle=\color{codecomment}\itshape,
    stringstyle=\color{codestring},
    aboveskip=1em,
    belowskip=1em,
    columns=fullflexible,
}

\lstdefinestyle{textstyle}{
    backgroundcolor=\color{codegray},
    basicstyle=\ttfamily\footnotesize,
    breaklines=true,
    frame=single,
    rulecolor=\color{codeborder},
    showstringspaces=false,
    aboveskip=1em,
    belowskip=1em,
    columns=fullflexible,
}

\lstset{style=textstyle}

% --- Definition / Idea / Important boxes ---
\usepackage[most]{tcolorbox}
\newtcolorbox{definicion}[1][]{
    colback=blue!5!white,
    colframe=blue!50!black,
    fonttitle=\bfseries,
    title=Definición,
    sharp corners=southwest,
    rounded corners=northwest,
    arc=2pt,
    boxrule=0.6pt,
    left=8pt,right=8pt,top=4pt,bottom=4pt,
    #1
}
\newtcolorbox{idea}[1][]{
    colback=orange!5!white,
    colframe=orange!60!black,
    fonttitle=\bfseries,
    title=Idea intuitiva,
    sharp corners=southwest,
    rounded corners=northwest,
    arc=2pt,
    boxrule=0.6pt,
    left=8pt,right=8pt,top=4pt,bottom=4pt,
    #1
}
\newtcolorbox{importante}[1][]{
    colback=red!5!white,
    colframe=red!50!black,
    fonttitle=\bfseries,
    title=Importante,
    sharp corners=southwest,
    rounded corners=northwest,
    arc=2pt,
    boxrule=0.6pt,
    left=8pt,right=8pt,top=4pt,bottom=4pt,
    #1
}

% --- Hyperlinks ---
\usepackage[colorlinks=true,
            linkcolor=blue!60!black,
            urlcolor=blue!50!black,
            citecolor=blue!60!black]{hyperref}

% --- Color palette ---
\definecolor{io1}{HTML}{4C72B0}
\definecolor{io2}{HTML}{DD8452}
\definecolor{io3}{HTML}{55A868}
\definecolor{io4}{HTML}{8172B2}
\definecolor{ioctl}{HTML}{C44E52}

% --- Title ---
\title{
    \vspace{-2.5em}
    \textbf{Resumen teórico --- Sistemas Operativos} \\[0.3em]
    \large Notas del curso: \emph{Entrada/Salida}
}
\author{
    Tomás Rodeghiero \\
    \small Sistemas Operativos --- Universidad Nacional de Río Cuarto
}
\date{2026}

\begin{document}

\maketitle

\begin{abstract}
\noindent
Este documento es un resumen teórico del capítulo \emph{Entrada/Salida}
de las \emph{Notas del curso} de Sistemas Operativos (Marcelo Arroyo,
UNRC). Se presentan la clasificación de los dispositivos, las formas de
interconectarlos a la CPU (bus de E/S dedicado y \emph{memory-mapped
I/O}), el modelo de control basado en puertos, la arquitectura de los
\emph{device drivers}, el diseño del subsistema de E/S, las disciplinas
de líneas, los buffers/caché de bloques, la planificación de
requerimientos (incluido el algoritmo del elevador), la transferencia
directa a memoria (DMA), la integración con el sistema de archivos
(archivos especiales en \texttt{/dev}, \emph{major/minor numbers},
\texttt{file\_operations}), los dispositivos virtuales (\texttt{random},
\texttt{zero}, \texttt{null}, RAID), los \emph{kernel modules}
cargables dinámicamente y los estándares más relevantes (IDE/ATA,
UART, VirtIO). El objetivo es servir como material de estudio para
preparar los prácticos asociados a E/S y el final de la materia.
\end{abstract}

\tableofcontents
\bigskip

% =====================================================================
\section{Introducción y clasificación de dispositivos}
% =====================================================================

Un sistema de computación incluye generalmente \emph{muchos} dispositivos
de entrada/salida (E/S), de distintos tipos. Hay dos formas naturales
de clasificarlos.

\subsection{Por dirección de transmisión de datos}

\begin{itemize}[leftmargin=*,nosep]
    \item \textbf{Entrada}: sólo se pueden leer datos del dispositivo
        (por ejemplo, un teclado, un mouse o un sensor).
    \item \textbf{Salida}: sólo se pueden escribir datos hacia el
        dispositivo (por ejemplo, una impresora o un display).
    \item \textbf{Entrada-Salida}: se pueden hacer ambas operaciones
        (un disco, una interfaz de red).
\end{itemize}

\subsection{Por modo de transferencia}

\begin{itemize}[leftmargin=*]
    \item \textbf{Dispositivos de caracteres (bytes).} Transfieren de a
        bytes. Ejemplos: teclado, mouse, dispositivos de
        comunicaciones tipo UART.
    \item \textbf{Dispositivos de bloques.} Transfieren de a
        \emph{bloques} de tamaño fijo (por ejemplo, 512\,bytes). El
        ejemplo paradigmático es el disco.
\end{itemize}

Hay dispositivos que no encajan limpiamente en esta clasificación: por
ejemplo, las interfaces de red tipo Ethernet transmiten en forma
serial bloques de tamaño \emph{variable} llamados \textbf{paquetes}
(\emph{packets}), que tienen un tamaño mínimo y uno máximo.

\paragraph{El timer.}
Otro dispositivo que se maneja de forma especial es el \textbf{timer}:
se programa para generar \emph{interrupciones periódicas} y así
permitirle al kernel tomar el control con regularidad (es la base de
la preemption del scheduler y del envejecimiento de bits en memoria
virtual).

% =====================================================================
\section{Interconexión de dispositivos}
% =====================================================================

Una arquitectura puede conectar sus dispositivos a la CPU de dos
maneras:

\begin{enumerate}[leftmargin=*]
    \item Mediante un \textbf{bus de E/S} dedicado y separado del bus
        de memoria.
    \item Conectándolos directamente al espacio de direcciones de
        memoria: \textbf{memory-mapped I/O (MMIO)}.
\end{enumerate}

La interconexión al espacio de direcciones puede ser \emph{hard-wired},
pero las arquitecturas modernas usan un \textbf{memory controller chip}
programable que permite rutear los requerimientos de lectura/escritura
a distintos chips (ROM, DRAM, dispositivos de I/O) según el valor del
bus de direcciones.

\subsection{Ejemplo: arquitectura x86 (híbrida)}

La arquitectura x86 (Intel) es \textbf{híbrida}: la CPU contiene un
\emph{bus de E/S} separado del bus de memoria. Algunos dispositivos se
conectan al bus de E/S (con un espacio de direcciones de 16\,bits),
mientras que otros se mapean en memoria.

\begin{itemize}[leftmargin=*,nosep]
    \item Para el bus de E/S, la CPU provee instrucciones especiales:
        \texttt{in} y \texttt{out}.
    \item En una PC basada en x86, el \emph{frame buffer} del
        controlador de video (CGA, VGA, etc.) está \textbf{mapeado en
        memoria}: al escribir datos en ese espacio reservado se produce
        una salida en el display.
\end{itemize}

\subsection{Ejemplo: RISC-V}

El board RISC-V usado en el taller no dispone de un \emph{video
controller}, pero soporta la conexión de una terminal (\emph{tty}) por
su puerto serie UART. La terminal transmite los códigos de las teclas
pulsadas y muestra en pantalla los caracteres recibidos. La interfaz
de programación del UART permite transmitir y recibir bytes mediante
un protocolo serial muy simple.

\begin{idea}
Las arquitecturas modernas como RISC-V tienden a usar exclusivamente
\textbf{MMIO}, aprovechando los enormes espacios de direcciones lógicas
disponibles. Esto simplifica el diseño del hardware y elimina la
necesidad de un bus adicional.
\end{idea}

\subsection{Ejemplo: modo texto VGA}

El controlador VGA en la IBM PC expone el \emph{text buffer} de video
en la dirección física \texttt{0xB800}, de tamaño dependiente del modo
(típicamente 80 columnas $\times$ 25 filas). Cada carácter se
representa con 16\,bits de la forma \emph{(ascii, attributes)}: el byte
ASCII más un byte de atributos (color de fondo, color de frente y
parpadeo). En DOS o en las primeras versiones de MS-Windows, las
aplicaciones podían escribir directamente en este buffer, sin
protección.

\paragraph{Modo gráfico y fonts.}
El hardware gráfico moderno soporta \textbf{modo gráfico} (matriz de
píxeles). El modo texto usado en consolas se implementa
\emph{dibujando} los caracteres gráficamente. La forma de cada carácter
se describe en un \textbf{font}, que puede ser \emph{raw} (matriz de
píxeles) o \emph{vectorial} (paths de trazado). La representación
vectorial permite un mejor escalado.

% =====================================================================
\section{Buses de interconexión}
% =====================================================================

La mayoría de las arquitecturas soportan la conexión de \emph{buses de
E/S}. Un bus es a su vez un dispositivo con su propio controlador
programable. Permite conectar varios dispositivos y, además, otros
buses secundarios.

\begin{center}
\begin{tikzpicture}[
    every node/.style={font=\ttfamily\small},
    box/.style={draw,minimum width=2.0cm,minimum height=0.7cm,align=center},
    line/.style={thick}
]
\node[box] (cpu) at (0,2) {cpu};

\node[box,fill=io1!15] (rom) at (3.5,3.6) {ROM};
\node[box,fill=io2!15,below=0pt of rom] (ram) {RAM};
\node[box,fill=io3!15,below=0pt of ram] (io) {I/O ports};

\node[box,fill=io4!20] (mmio) at (8.0,2.2) {memory\\mapped\\devices};

\draw[line] (cpu) -- node[above,font=\scriptsize]{mem bus} ($(rom.west)!0.5!(io.west)$);
\draw[<->,line] (io.east) -- (mmio.west);

% I/O bus and PCI
\node[box] (dev) at (3.5,-1.0) {device};
\node[box] (usbhub) at (6.5,-1.0) {USB hub};
\node[box] (usbdev) at (9.5,-1.0) {usb device};

\draw[line] (cpu.south) -- ++(0,-0.8) -| node[pos=0.25,above,font=\scriptsize]{I/O bus} (dev.north);
\draw[line] (cpu.south) -- ++(0,-1.4) -| node[pos=0.30,above,font=\scriptsize]{PCI bus} (usbhub.north);
\draw[line] (usbhub.east) -- (usbdev.west);
\end{tikzpicture}
\end{center}

\subsection{Buses estándar: PCI y USB}

Un bus estándar muy usado como bus principal es el \textbf{PCI}
(\emph{Peripheral Component Interconnect}). Permite conectar:

\begin{itemize}[leftmargin=*,nosep]
    \item interfaces de red,
    \item controladoras de disco,
    \item co-procesadores gráficos,
    \item dispositivos de audio,
    \item y a su vez otros controladores como \textbf{USB}
        (\emph{Universal Serial Bus}), que soporta un árbol de
        \emph{USB hubs}, donde cada hub contiene puertos en los que se
        conectan dispositivos físicos.
\end{itemize}

Los estándares de buses como PCI o USB describen tanto las
características \emph{físicas} (conectores, niveles eléctricos) como la
\emph{interfaz de control} (registros de control, estado y datos, y
los detalles de programación).

\subsection{Auto-detección y plug-and-play}

Cada dispositivo compatible \textbf{se anuncia en el bus} al iniciarse
(o al ser conectado), informando: tipo de dispositivo, fabricante,
registros de E/S, líneas de interrupción a usar, etc. El controlador
del bus registra esta información en una tabla que permite al SO
\emph{auto-detectar} los dispositivos existentes.

Algunos buses permiten conexión en caliente: notifican al kernel
(generando una interrupción) cuando se conecta un nuevo dispositivo,
permitiendo activar el \emph{device driver} correspondiente. Esto se
conoce como \textbf{plug-and-play}.

% =====================================================================
\section{Control de dispositivos}
% =====================================================================

Un dispositivo expone un conjunto de \textbf{puertos} (o registros) de
control, estado y datos. Estos puertos están contenidos en el
\textbf{controlador del dispositivo} (\emph{device controller}).

\begin{center}
\begin{tikzpicture}[
    every node/.style={font=\ttfamily\small},
    reg/.style={draw,minimum width=2.4cm,minimum height=0.8cm,align=center}
]
\node[font=\bfseries\small] at (0,3.0) {BASE\_ADDR};
\node[reg,fill=io1!15] (c) at (0,2.3) {control};
\node[reg,fill=io2!15,below=0pt of c] (s) {status};
\node[reg,fill=io3!15,below=0pt of s,minimum height=1.3cm] (d) {data};
\end{tikzpicture}
\end{center}

\begin{itemize}[leftmargin=*]
    \item \textbf{Registros de control.} Se usan para enviar comandos
        al dispositivo.
    \item \textbf{Registros de datos.} Contienen los datos a transferir
        (lectura o escritura).
    \item \textbf{Registros de estado.} Indican el estado del
        dispositivo: si hay un dato disponible, si se produjo un error
        o si el dispositivo está configurado para generar interrupciones.
\end{itemize}

Estos puertos aparecen en el espacio de direcciones del bus de E/S o
en el de memoria, según la forma de interconexión. La documentación
técnica del dispositivo (o del estándar si lo cumple) describe cómo
debe el SO interactuar con él; las direcciones suelen ser
\emph{dependientes de la plataforma}.

% =====================================================================
\section{Device drivers}
% =====================================================================

\begin{definicion}
Un \textbf{device driver} es el módulo del kernel que controla un tipo
o familia de dispositivos. El SO define la \emph{interfaz} (funciones)
que el driver debe implementar; el kernel se comunica con el driver
mediante esa API.
\end{definicion}

\paragraph{Diseño polimórfico.}
Se sigue un diseño polimórfico al estilo \emph{orientado a objetos}:
cada device driver expone una estructura de datos con \emph{punteros}
a las funciones que implementan la API. El kernel invoca esos punteros
sin conocer los detalles internos del driver.

\subsection{Funciones que debe implementar un driver}

Comúnmente un device driver implementa:

\begin{itemize}[leftmargin=*]
    \item \textbf{Detección} (\emph{discovery}) de los dispositivos que
        puede controlar.
    \item \textbf{Inicialización} de los dispositivos: configurar las
        interrupciones que generará, informar al kernel las direcciones
        de E/S que reservará, etc.
    \item \textbf{Operaciones de lectura y escritura}: implementar el
        control y programación de los requerimientos de transferencia
        (\emph{I/O requests}).
    \item \textbf{Manejo de interrupciones}: cuando un dispositivo
        completa una operación, genera una interrupción y el driver
        debe atenderla.
\end{itemize}

\subsection{Polling como alternativa}

Una forma alternativa de detectar la finalización de una operación es
\textbf{pooling} (o \emph{polling}): verificar periódicamente un bit
de un registro de estado del dispositivo.

\begin{importante}
Las interrupciones son preferibles cuando el dispositivo es ``lento''
(el CPU puede hacer otra cosa mientras espera). El polling es preferible
cuando el dispositivo es ``rápido'' y se espera un dato inmediato; el
costo de manejar la interrupción superaría al de simplemente leer el
estado.
\end{importante}

% =====================================================================
\section{Diseño del subsistema de E/S}
% =====================================================================

Dada la enorme variedad de dispositivos, el SO necesita un diseño que
abstraiga a los device drivers de las capas superiores del kernel.

Cada device driver se \textbf{registra} en el sistema y reserva los
recursos que requiere: espacios de direcciones para los puertos del
controlador, y la línea (o número) de interrupción que usará. El
subsistema de E/S mantiene entonces:

\begin{itemize}[leftmargin=*,nosep]
    \item Una \textbf{tabla de dispositivos}, cuya estructura por
        entrada contiene las direcciones de puertos de I/O, las
        líneas de IRQ y otros datos.
    \item Una asociación entre cada dispositivo y su \emph{device
        driver}, representado por una estructura de datos con punteros
        a sus operaciones soportadas.
\end{itemize}

\subsection{Interfaz de un device driver}

Para facilitar la interoperabilidad, cada driver implementa una interfaz
definida en el SO (habitualmente por el filesystem). Es común que se
defina una interfaz \emph{por cada tipo de dispositivo}. En los SO tipo
UNIX hay tres clases:

\begin{enumerate}[leftmargin=*]
    \item \textbf{Dispositivos de caracteres.} La E/S es de a bytes.
        Ejemplo: teclado, mouse, puerto serie.
    \item \textbf{Dispositivos de bloques.} La lectura/escritura es por
        bloques (discos, almacenamiento no volátil). El filesystem
        ofrece una caché para la E/S de bloques.
    \item \textbf{Dispositivos de comunicaciones} o \textbf{red}.
        Ejemplo: interfaces Ethernet, Wifi, Bluetooth.
\end{enumerate}

% =====================================================================
\section{Software independiente de dispositivos}
% =====================================================================

El subsistema de E/S provee además algunos servicios o políticas que
\emph{no dependen} del dispositivo concreto: planificación de
requerimientos, manejo de buffers, etc.

\subsection{Disciplina de líneas (line discipline)}

Un ejemplo importante es el servicio de \textbf{disciplina de líneas}
entre los \emph{tty device drivers} y las llamadas al sistema como
\texttt{read}/\texttt{write} sobre terminales de texto. Una disciplina
típica realiza E/S de texto \emph{por líneas} (cadenas de caracteres).

\paragraph{Entrada.}
La disciplina de entrada (\emph{input}) tiene en cuenta caracteres de
control:

\begin{itemize}[leftmargin=*,nosep]
    \item \texttt{Enter}: indica fin de la línea.
    \item \texttt{Ctrl-D}: indica fin de la entrada (EOF).
    \item \texttt{BS} (\emph{backspace}): borra el carácter ingresado
        previamente.
\end{itemize}

Su implementación se basa en la gestión de un \textbf{buffer de
caracteres de entrada organizado como una cola circular}.

\paragraph{Salida.}
En la salida (\emph{output}) se interpretan caracteres de control como:

\begin{itemize}[leftmargin=*,nosep]
    \item \texttt{CR} (\emph{carriage return}): retorno de carro.
    \item \texttt{LF} (\emph{linefeed}): bajada de cursor (scroll) y
        avance de línea.
\end{itemize}

\begin{idea}
En xv6, el módulo \texttt{uart.c} implementa el driver y la disciplina
de líneas sobre una terminal conectada al puerto serie UART.
\end{idea}

% =====================================================================
\section{Buffers y caché}
% =====================================================================

Para minimizar los accesos al dispositivo (como un disco), cada vez que
se lee un bloque se carga en bloques de memoria (típicamente páginas)
que se manejan como una \textbf{caché}. Esa caché tiene un límite de
bloques.

\paragraph{Política de reemplazo.}
Los bloques de la caché se \emph{reúsan}; comúnmente se usa un
algoritmo \textbf{LRU} para elegir la víctima a liberar (estamos viendo
la misma idea que en memoria virtual, ahora aplicada a bloques de
disco).

\paragraph{Sincronización a disco.}
Las llamadas al sistema \verb|flush(fd)| o \verb|close(fd)| producen
que se transfieran los bloques modificados (\emph{dirty}) del archivo
desde la caché al disco.

\begin{importante}
La E/S de dispositivos de \emph{caracteres} comúnmente \textbf{no}
utiliza buffers de caché: los datos se procesan a medida que llegan
(por ejemplo, una pulsación de tecla).
\end{importante}

% =====================================================================
\section{Planificación de E/S}
% =====================================================================

El SO debe gestionar los requerimientos de E/S generados por los
procesos (y por el propio kernel). Algunos algoritmos son similares a
los de planificación de procesos:

\begin{itemize}[leftmargin=*,nosep]
    \item \textbf{FIFO}: los requerimientos se procesan en orden de
        llegada.
    \item \textbf{Prioridades}: se ordenan según una prioridad asignada.
\end{itemize}

\subsection{Discos: algoritmo del elevador}

Para dispositivos de bloques electromecánicos (discos magnéticos), los
requerimientos pueden ordenarse por la \textbf{cercanía a la posición
actual del brazo} de lectura/escritura. Procesar los más cercanos
reduce significativamente los tiempos de acceso. Pero esto puede
hacer que requerimientos en cilindros lejanos esperen indefinidamente
(\emph{starvation}). La solución es el \textbf{algoritmo del elevador}:

\begin{definicion}
El \textbf{algoritmo del elevador} mantiene una variable
\texttt{direction} que representa la dirección actual del brazo. Los
requerimientos se ordenan y procesan en \emph{esa} dirección. Al
llegar al cilindro del extremo, se invierte \texttt{direction} y se
procesan los restantes en orden de cercanía en la nueva dirección.
\end{definicion}

Es análogo al funcionamiento de un ascensor: atiende a todos los pisos
en una dirección antes de cambiar de sentido.

% =====================================================================
\section{Transferencia directa a memoria (DMA)}
% =====================================================================

Hay dos formas de mover datos entre la RAM y los buffers internos de un
dispositivo:

\begin{enumerate}[leftmargin=*]
    \item \textbf{Por software}: el driver hace una copia byte a byte
        o palabra a palabra, del estilo de un \texttt{memcpy}
        (\texttt{dst}, \texttt{src}, \texttt{count}). Esto puede
        insumir muchos ciclos de CPU.
    \item \textbf{Por DMA} (\emph{Direct Memory Access}): el dispositivo
        realiza la transferencia de manera autónoma, en paralelo con la
        ejecución de instrucciones de la CPU.
\end{enumerate}

\subsection{Cómo funciona el DMA}

Para usar DMA, el device driver le indica al dispositivo:

\begin{enumerate}[leftmargin=*,nosep]
    \item La \textbf{dirección física} del buffer en memoria.
    \item El \textbf{tamaño} de la transferencia.
    \item La orden de iniciar.
\end{enumerate}

Mientras la transferencia se realiza, el dispositivo DMA compite con la
CPU por el acceso al bus de memoria mediante una técnica de
\textbf{robo de ciclos} (\emph{cycle stealing}): aprovecha los ciclos
en los que la CPU no necesita la memoria.

\paragraph{Notificación.}
Cuando la transferencia finaliza, el dispositivo DMA \textbf{genera una
interrupción} para avisarle al driver.

\begin{idea}
Actualmente, muchos dispositivos modernos ya incluyen un mecanismo DMA
\emph{en el mismo chip}, por lo que el driver no necesita programar un
controlador DMA externo.
\end{idea}

% =====================================================================
\section{Integración con el sistema de archivos}
% =====================================================================

El sistema de E/S se integra con el sistema de archivos para brindar
una abstracción central en los SO tipo UNIX:

\begin{importante}
Los dispositivos \textbf{son archivos especiales}, y los procesos
acceden a ellos (a través de los drivers) mediante la API de
operaciones sobre archivos: \texttt{open}, \texttt{read},
\texttt{write}, \texttt{close}, \ldots
\end{importante}

\subsection{La estructura \texttt{file\_operations} en Linux}

El filesystem define un \emph{conjunto de operaciones} que el driver de
ese tipo debe implementar. En Linux, esta interfaz toma la forma de
una \texttt{struct} con punteros a función:

\begin{lstlisting}[style=cstyle]
struct file_operations {
    struct module *owner;
    loff_t   (*llseek)  (struct file *, loff_t, int);
    ssize_t  (*read)    (struct file *, char *, size_t, loff_t *);
    ssize_t  (*write)   (struct file *, const char *, size_t, loff_t *);
    int      (*readdir) (struct file *, void *, filldir_t);
    unsigned (*poll)    (struct file *, struct poll_table_struct *);
    int      (*ioctl)   (struct inode *, struct file *, unsigned, unsigned long);
    int      (*mmap)    (struct file *, struct vm_area_struct *);
    int      (*open)    (struct inode *, struct file *);
    int      (*flush)   (struct file *);
    int      (*release) (struct inode *, struct file *);
    int      (*fsync)   (struct file *, struct dentry *, int);
    int      (*fasync)  (int, struct file *, int);
    int      (*lock)    (struct file *, int, struct file_lock *);
    ssize_t  (*readv)   (struct file *, const struct iovec *, ...);
    ssize_t  (*writev)  (struct file *, const struct iovec *, ...);
};
\end{lstlisting}

Cada driver llena esta estructura con sus implementaciones específicas.

\subsection{Archivos especiales en \texttt{/dev}}

En los sistemas tipo UNIX, los archivos que representan dispositivos
están creados bajo \texttt{/dev}. El \emph{tipo} de cada archivo
indica la clase de dispositivo:

\begin{itemize}[leftmargin=*,nosep]
    \item \texttt{c} si es de caracteres,
    \item \texttt{b} si es de bloques.
\end{itemize}

Un proceso de usuario puede usar \texttt{open}, \texttt{read},
\texttt{write}, \texttt{close} sobre estos archivos para interactuar
directamente con el driver del kernel. Por ejemplo:

\begin{lstlisting}[style=shellstyle]
$ echo hello > /dev/tty
hello
\end{lstlisting}

\subsection{Major y minor numbers}

En los sistemas tipo UNIX, un archivo especial contiene sólo
\emph{metadatos} (no almacena datos en disco). Los metadatos más
importantes son:

\begin{itemize}[leftmargin=*]
    \item Tipo (\texttt{c} o \texttt{b}).
    \item \textbf{Major number}: representa el \emph{tipo} de
        dispositivo y le permite al kernel identificar al device
        driver correspondiente.
    \item \textbf{Minor number}: representa el \emph{número de
        dispositivo} dentro de los de su tipo. Se pasa al driver en
        cada operación (en un campo de la estructura \texttt{inode}),
        para que éste sepa con qué dispositivo concreto interactuar.
\end{itemize}

\begin{idea}
Esta separación permite que un \emph{único driver} maneje muchos
dispositivos del mismo tipo: el major identifica al driver y el minor
le dice ``cuál de mis dispositivos''.
\end{idea}

% =====================================================================
\section{Dispositivos virtuales}
% =====================================================================

Es común que el SO implemente dispositivos en \emph{software} (no
existen físicamente):

\begin{itemize}[leftmargin=*]
    \item \texttt{random}: generador de números aleatorios. Sólo
        lectura.
    \item \texttt{zero}: retorna ceros. Sólo lectura.
    \item \texttt{null}: ignora todo lo que recibe. Sólo escritura.
\end{itemize}

Otros dispositivos virtuales agrupan varios dispositivos físicos y usan
sus drivers correspondientes:

\paragraph{Consola/terminal (tty).}
Una \emph{consola} consiste en un teclado y un display. El driver de la
\emph{tty} define el control y la asociación con esos dispositivos. Un
driver de este tipo suele soportar \textbf{consolas o terminales
virtuales}: el SO permite al usuario crear varias terminales como
espacios de trabajo distintos.

\paragraph{RAID.}
Un \textbf{RAID} (\emph{Redundant Array of Inexpensive Disks}) es un
conjunto de discos presentado al sistema como un único disco. Es un
dispositivo de almacenamiento tolerante a fallas que permite accesos
en paralelo:

\begin{itemize}[leftmargin=*,nosep]
    \item \textbf{Distribuye} los bloques de datos entre los discos:
        un archivo cuyos bloques están dispersos puede leerse en
        paralelo.
    \item \textbf{Replica} los bloques: ante la falla de un disco,
        existe una copia en otro y el sistema sigue funcionando.
\end{itemize}

Existe hardware específico para RAID, pero también puede implementarse
por software usando discos estándar (por ejemplo, ATA). GNU/Linux
soporta RAID por software.

% =====================================================================
\section{Device drivers cargables dinámicamente}
% =====================================================================

La mayoría de los SO modernos permiten cargar device drivers
\textbf{dinámicamente}, es decir, sin necesidad de reiniciar ni
recompilar el kernel. Esto significa que el driver se compila como un
\textbf{kernel module}.

\subsection{Diseño microkernel vs.\ monolítico}

\begin{itemize}[leftmargin=*]
    \item En un diseño \textbf{microkernel}, esto es natural: cada
        driver es un componente (programa) independiente del resto del
        kernel y se comunica con el kernel \emph{por mensajes}. Por
        ejemplo, en MS-Windows los drivers toman la forma de
        \emph{Dynamic Link Libraries} (DLLs).
    \item En un diseño \textbf{monolítico}, es más complicado: el
        módulo debe enlazarse dinámicamente con el código del kernel
        para vivir en su \emph{mismo espacio de memoria}.
\end{itemize}

\subsection{Linux Kernel Modules (LKM)}

En GNU/Linux los módulos son \textbf{shared objects} compilados y
enlazados como una biblioteca compartida con extensión \texttt{.ko}
(\emph{kernel object}). Se llaman \textbf{Linux Kernel Modules (LKM)}.

\begin{itemize}[leftmargin=*]
    \item El comando \texttt{insmod} mapea el archivo objeto (driver)
        en memoria y ejecuta la llamada al sistema
        \verb|init_module()|, que realiza el enlazado dinámico del
        módulo en el espacio del kernel e invoca a su función
        \texttt{init()}.
    \item GNU/Linux provee los comandos de usuario \texttt{lsmod},
        \texttt{insmod} y \texttt{rmmod} para listar, cargar y remover
        módulos cargables dinámicamente.
\end{itemize}

% =====================================================================
\section{Estándares de dispositivos}
% =====================================================================

Existen varios estándares para diferentes tipos de dispositivos. Estos
estándares especifican detalles \emph{físicos} (niveles eléctricos,
conectores) y \emph{lógicos} (registros de control/estado, detalles de
configuración y operación).

\begin{idea}
Que los dispositivos adhieran a estándares simplifica drásticamente la
escritura de los drivers: basta con basarse en la documentación del
estándar, sin tener que adaptarse a cada fabricante.
\end{idea}

Algunos estándares relevantes:

\begin{itemize}[leftmargin=*]
    \item \textbf{IDE} (\emph{Integrated Drive Electronics}): interfaz
        estándar para dispositivos de almacenamiento masivo (discos).
    \item \textbf{AT Attachment (ATA)}: evolución de IDE; define
        dispositivos con transferencia en paralelo (\textbf{PATA}) y
        en serie (\textbf{SATA}).
    \item \textbf{UART} (\emph{Universal Asynchronous
        Receiver-Transmitter}): dispositivos de comunicación serial.
    \item \textbf{ISO/IEC JTC}: estándares para equipos de oficina
        como impresoras y escáneres.
    \item \textbf{VirtIO}: interfaz de dispositivos \emph{virtuales}.
        Pensada para SO corriendo sobre máquinas virtuales. Define
        diferentes tipos de dispositivos y sus interfaces \emph{lógicas}
        (puertos y operación). \emph{No} incluye detalles físicos, ya
        que su objetivo es definir dispositivos virtuales emulados por
        software.
\end{itemize}

% =====================================================================
\section{VirtualIO interface}
% =====================================================================

VirtIO define una \textbf{interfaz unificada} para diferentes tipos de
dispositivos. Los dispositivos VirtIO aparecen mapeados en memoria o
conectados a un bus PCI y exponen:

\begin{itemize}[leftmargin=*,nosep]
    \item su \emph{tipo} (block, network, input device, \ldots),
    \item su \emph{estado},
    \item las \emph{características} que soportan,
    \item el mecanismo de \emph{notificaciones},
    \item un espacio de \emph{configuración},
    \item y una o más \textbf{virtual queues}.
\end{itemize}

\subsection{Virtual queues}

\begin{definicion}
Una \textbf{virtual queue} es una estructura que permite gestionar
buffers de entrada/salida \emph{compartidos} entre el device driver y
el device controller (el dispositivo). Es el canal de transferencia de
datos y requerimientos entre ambos.
\end{definicion}

\paragraph{Uso típico.}

\begin{itemize}[leftmargin=*]
    \item Con dispositivos como \textbf{discos}, el driver usa
        \emph{una} virtual queue con los requerimientos y buffers de
        lectura/escritura.
    \item Con dispositivos de \textbf{red}, se usan \emph{dos} virtual
        queues: una para los paquetes recibidos y otra para los
        paquetes a enviar.
\end{itemize}

% =====================================================================
\section{Cierre}
% =====================================================================

El subsistema de E/S es la capa que pone en contacto al kernel con la
diversidad infinita de hardware que existe afuera. La forma en que
las Notas organizan esta diversidad puede resumirse en pocas ideas:

\begin{itemize}[leftmargin=*]
    \item \textbf{Clasificar} los dispositivos por dirección de
        transmisión y por modo de transferencia (caracteres, bloques,
        paquetes) permite definir interfaces uniformes.
    \item \textbf{Conectar} los dispositivos a la CPU vía buses
        dedicados o vía \emph{memory-mapped I/O}; las arquitecturas
        modernas (RISC-V) se inclinan a MMIO porque aprovecha el
        espacio de direcciones lógicas y simplifica el hardware.
    \item \textbf{Programar} cada dispositivo mediante un conjunto de
        \emph{puertos} de control, datos y estado, encapsulados en su
        controlador.
    \item \textbf{Abstraer} cada dispositivo con un \textbf{device
        driver} que implementa una API polimórfica común (detección,
        init, read/write, manejo de IRQ). El subsistema de E/S
        mantiene la tabla de dispositivos y la asociación driver
        $\leftrightarrow$ dispositivo.
    \item \textbf{Compartir} código entre drivers mediante \emph{software
        independiente del dispositivo}: disciplinas de líneas, buffer
        de caché de bloques (con LRU), planificación de requerimientos
        (FIFO, prioridades, algoritmo del elevador en discos).
    \item \textbf{Acelerar} las transferencias mediante \textbf{DMA},
        que libera a la CPU del trabajo de copiar datos por software.
    \item \textbf{Integrar} todo en el filesystem: los dispositivos son
        archivos especiales en \texttt{/dev}, con \emph{major}/\emph{minor
        numbers}, accesibles vía la API estándar \texttt{open}/
        \texttt{read}/\texttt{write}/\texttt{close}; en Linux la
        \texttt{struct file\_operations} concreta esta idea.
    \item \textbf{Extender} el sistema con dispositivos virtuales
        (\texttt{random}, \texttt{zero}, \texttt{null}, tty virtuales,
        RAID por software) y con módulos cargables dinámicamente
        (LKM, \texttt{insmod}, \texttt{rmmod}).
    \item \textbf{Estandarizar} las interfaces de hardware (IDE/ATA,
        UART, VirtIO) para simplificar la escritura de drivers; VirtIO
        en particular hace lo mismo pero para dispositivos virtuales
        en máquinas virtuales, usando \emph{virtual queues}.
\end{itemize}

\begin{idea}
La gran idea es la \textbf{uniformidad}: todo lo que pasa en E/S, por
distinto que sea el dispositivo, termina presentándose al resto del
kernel y al usuario como ``otro archivo''. Ese principio Unix --- "todo
es un archivo" --- es lo que hace manejable la complejidad inmensa del
hardware moderno.
\end{idea}

\end{document}
